%==============================================================================
\chapter{Design (UML models)}
\label{chap:conception}
%==============================================================================
The design file comprises the following UML diagrams. Each one is described by
its expected content in order to guide the modelling.
\begin{xltabular}{\textwidth}{>{\bfseries}p{3.8cm} X}
\toprule
\textbf{Diagram} & \textbf{Expected content} \\
\midrule
\endhead
Use case & Actors (farmer, agent, supervisor, manager, administrator, third-party application) and main cases (see figure \ref{fig:uc}). \\
Class & Entities of the data dictionary and their relationships (Farmer 1 to n Farm, Farm 1 to n Site, Site 1 to n Hive, Hive 1 to n Body or super, Body or super 1 to n Frame) and service methods. \\
Object & Snapshot of a site containing three populated hives with their frames. \\
Sequence & Course of carrying out a visit (Agent, React client, REST controller, business service, repository, database). \\
Activity & Process of planning, approving, executing and closing a visit. \\
State machine & Life cycle of a hive (created, populated, active, splitting, harvesting, closed). \\
Interaction or collaboration & Same scenario as the sequence, in the form of a collaboration between objects. \\
Package & Presentation, control, business, data-access, internationalisation, configuration and web-services layers. \\
Component & Application components (web application, OAuth service, API service, i18n module, configuration module, database). \\
Deployment & Distribution across the browser, the reverse proxy, the application container, the DBMS, the Keycloak identity provider and the sensors. \\
\bottomrule
\end{xltabular}

\section{Class diagram}
The class diagram, presented in horizontal alignment, describes the business
entities, their attributes, the service method \texttt{getZummHoneyActualQuantity}
and the relationships with their cardinalities.
\begin{figure}[htbp]
\centering
\includegraphics[width=\linewidth]{diagramme_classes}
\caption{Class diagram of the Zümm system (horizontal alignment).}
\label{fig:classes}
\end{figure}

\section{Layered view (package)}
\begin{figure}[htbp]
\centering
\begin{tikzpicture}[
  couche/.style={draw=mielfonce,fill=grisdoux,thick,minimum width=10cm,minimum height=8mm,font=\small}
]
\node[couche] (p) at (0,4) {Presentation: React and TypeScript};
\node[couche] (c) at (0,3) {Control: Spring MVC (REST API)};
\node[couche] (m) at (0,2) {Business: Spring services};
\node[couche] (d) at (0,1) {Persistence: Spring Data JPA and jOOQ};
\node[couche] (b) at (0,0) {Relational database};
\node[draw=mielfonce,thick,fill=miel!15,minimum height=8mm,font=\footnotesize,align=center,text width=3.1cm,right=4mm of c] (s) {i18n XML\\ConfigZumm.ini\\Web services};
\draw[-{Latex[length=2mm]},mielfonce] (p)--(c);
\draw[-{Latex[length=2mm]},mielfonce] (c)--(m);
\draw[-{Latex[length=2mm]},mielfonce] (m)--(d);
\draw[-{Latex[length=2mm]},mielfonce] (d)--(b);
\end{tikzpicture}
\caption{Multi-tier architecture in loosely coupled layers.}
\label{fig:couches}
\end{figure}

\section{Deployment view}
\begin{figure}[htbp]
\centering
\begin{tikzpicture}[
  noeud/.style={draw=mielfonce,fill=grisdoux,thick,minimum width=3cm,minimum height=1.1cm,font=\footnotesize,align=center},
  fl/.style={-{Latex[length=2mm]},mielfonce,thick}
]
\node[noeud] (nav) {Browser\\(web client)};
\node[noeud,right=1.6cm of nav] (srv) {Zümm container\\(Spring Boot)};
\node[noeud,right=1.6cm of srv] (bd) {PostgreSQL\\PostGIS + TimescaleDB};
\node[noeud,below=1cm of srv] (oauth) {Keycloak\\(Google federation)};
\node[noeud,below=1cm of bd] (capt) {Sensors\\(automated part)};
\draw[fl] (nav)--(srv) node[midway,above,font=\scriptsize]{HTTPS};
\draw[fl] (srv)--(bd) node[midway,above,font=\scriptsize]{JDBC / pool};
\draw[fl] (srv)--(oauth) node[midway,right,font=\scriptsize]{OIDC};
\draw[fl] (capt)--(srv);
\end{tikzpicture}
\caption{Deployment diagram of the system.}
\label{fig:deploiement}
\end{figure}

\section{Sequence diagram (UC2)}
The sequence diagram below details use case UC2, carry out a visit and fill in
the report. It shows the passage through the layers (agent, React client, REST controller, business service,
database) as well as the online and offline cases with deferred
synchronisation.
\begin{figure}[htbp]
\centering
\includegraphics[width=0.92\linewidth]{diagramme_sequence_uc2}
\caption{Sequence diagram of use case UC2 (carry out a visit).}
\label{fig:sequence-uc2}
\end{figure}

\section{Authentication sequence (Google OAuth)}
The authentication flow follows the \emph{Authorization Code} flow of OAuth~2.0.
The browser never carries the client secret: the exchange of the authorization
code for the access token is done \textbf{server to server}, over an SSL channel
authenticated by an X.509 certificate. The application session is then opened by
means of a secure cookie (\texttt{HttpOnly}, \texttt{Secure}), and the profile
returned by Google is associated with an existing \texttt{Agent} (or creates one
pending authorisation). In the demonstration, a local authentication fallback is
provided in case the provider is unavailable.
\begin{figure}[htbp]
\centering
\begin{tikzpicture}[font=\scriptsize,
  msg/.style={-{Latex[length=1.6mm]},mielfonce,thick},
  rmsg/.style={-{Latex[length=1.6mm]},mielfonce,thick,dashed}
]
% ---- lifelines ----
\foreach \x/\name in {0/{Agent (browser)}, 6/{Zümm server}, 12/{Google OAuth}}{
  \node[draw=mielfonce,fill=grisdoux,thick,align=center,minimum height=7mm,inner xsep=6pt] at (\x,0) {\name};
  \draw[mielfonce,dashed] (\x,-0.45) -- (\x,-9.9);
}
% ---- messages ----
\draw[msg]  (0,-1.0) -- (6,-1.0)  node[midway,above]{1. Access to a protected page};
\draw[rmsg] (6,-1.7) -- (0,-1.7)  node[midway,above]{2. 302: redirection (client\_id, scope, state)};
\draw[msg]  (0,-2.4) -- (12,-2.4) node[midway,above]{3. Authentication + consent};
\draw[rmsg] (12,-3.1) -- (0,-3.1) node[midway,above]{4. 302: redirection + code + state};
\draw[msg]  (0,-3.8) -- (6,-3.8)  node[midway,above]{5. GET /callback (code, state)};
\draw[msg]  (6,-4.5) -- (12,-4.5) node[midway,above]{6. POST /token (code, secret) --- SSL / X.509};
\draw[rmsg] (12,-5.2) -- (6,-5.2) node[midway,above]{7. access\_token + id\_token (JWT)};
\draw[msg]  (6,-5.9) -- (12,-5.9) node[midway,above]{8. GET /userinfo (access\_token)};
\draw[rmsg] (12,-6.6) -- (6,-6.6) node[midway,above]{9. Profile (email, name)};
\node[draw=mielfonce,fill=miel!12,thick,align=center,inner sep=3pt] at (6,-7.5)
  {10. Verify \texttt{state}, associate the Agent,\\open the session (RBAC)};
\draw[rmsg] (6,-8.6) -- (0,-8.6) node[midway,above]{11. 302: application + secure session cookie};
\end{tikzpicture}
\caption{OAuth~2.0 authentication sequence (Authorization Code) with Google.}
\label{fig:seq-oauth}
\end{figure}

\section{Object diagram}
The object diagram presents a snapshot of a site hosting three hives, with the
detail of the composition of a hive into body, super and frames.
\begin{figure}[htbp]
\centering
\includegraphics[width=0.7\linewidth]{diagramme_objets}
\caption{Object diagram, snapshot of an apiary site.}
\label{fig:objets}
\end{figure}

\section{Activity diagram}
The activity diagram models the process of planning, approving and carrying out
a visit, with the online and offline variants.
\begin{figure}[htbp]
\centering
\includegraphics[width=0.7\linewidth]{diagramme_activite}
\caption{Activity diagram of the visit process.}
\label{fig:activite}
\end{figure}

\section{State-machine diagram}
The state-machine diagram describes the life cycle of a hive, from its creation
to its closure.
\begin{figure}[htbp]
\centering
\includegraphics[width=0.85\linewidth]{diagramme_etat}
\caption{State machine of the life cycle of a hive.}
\label{fig:etat}
\end{figure}

\section{Collaboration diagram}
The collaboration diagram takes up the UC2 scenario in the form of numbered
messages exchanged between the objects.
\begin{figure}[htbp]
\centering
\includegraphics[width=\linewidth]{diagramme_collaboration}
\caption{Collaboration diagram of use case UC2.}
\label{fig:collaboration}
\end{figure}

\section{Component diagram}
The component diagram shows the organisation of the application components, their
interfaces and their dependencies, including the service interface exposing the
honey quantity.
\begin{figure}[htbp]
\centering
\includegraphics[width=\linewidth]{diagramme_composant}
\caption{Component diagram of the Zümm system.}
\label{fig:composant}
\end{figure}
